% !TeX root = ../教材题库_全部条目_转写工作稿.tex
% Source: XJ-B1-C02-HW-P049-Y0793
% 第2章 / 2.2 从函数观点看一元二次方程 / 教材 p.49 / PDF 第56页
% OCR未校对：公式与文字均可能有误；未完成数学TeX转写；禁止直接发给学生
% 原题及相邻原解截图（核对用）：../题目原图/XJ-B1-C02-HW-P049-Y0793-01.webp
% 本文件只在转写完成后作为题干使用；OCR 条目尚不是可靠题干。
\begin{bankitem}{习题2.2 第4题}{第2章 / 2.2 从函数观点看一元二次方程 / 教材 p.49 / PDF 第56页}{XJ-B1-C02-HW-P049-Y0793}{OCR未校对：公式与文字均可能有误；未完成数学TeX转写；禁止直接发给学生}
\begin{ocrdraft}
[教材p.49 / PDF第56页]
\par
4．如图所示为二次函数3，一依r2十的图象，若一元二
\par
次方程依r2十十m一0有实数根，求m的最大值．
\end{ocrdraft}
\end{bankitem}
